\chapter{Aprendizaje de máquina: predecir}
\label{cap:ml}
\textit{Criterio 4 de la rúbrica (12\,\%): construye, valida, compara e interpreta modelos con métricas adecuadas y
justifica la selección.}

El proyecto tiene dos problemas de aprendizaje de máquina, y es importante no confundirlos:
\begin{center}
\begin{tabularx}{\linewidth}{lYY}
\encabezado & \h{Clasificación (Radar)} & \h{Regresión de series de tiempo (Pronóstico)} \\
\textbf{Qué predice} & Un \emph{estado}: tranquilo, concurrido o saturado & Un \emph{número}: cuántos visitantes \\
\textbf{Horizonte} & El mes siguiente & 1 a 12 meses \\
\textbf{Modelos} & Logística, Random Forest, Gradient Boosting & Línea base, Holt-Winters (2), regresión con clima, Gradient Boosting \\
\textbf{Contra qué se compara} & ``Igual que este mes'' & ``El mismo mes del año pasado'' \\
\textbf{Métrica} & Aciertos, cambios anticipados, F1 & MAE, MAPE, cobertura del rango \\
\bottomrule
\end{tabularx}
\end{center}

\section{Clasificación: ¿qué estado tendrá cada lugar el mes siguiente?}
\subsection{El planteamiento}
\textbf{Unidad}: un lugar en un mes. \textbf{Objetivo} ($y$): el estado del índice comparable el mes siguiente.
\textbf{Rasgos} ($x$, lo que el modelo ve):
\begin{itemize}
\item el índice de este mes y de los dos anteriores ($IPT_t$, $IPT_{t-1}$, $IPT_{t-2}$) y el cambio ($IPT_t-IPT_{t-1}$);
\item el mes del año como un punto en un círculo, $\left(\sin\frac{2\pi m}{12},\cos\frac{2\pi m}{12}\right)$. \emph{Por qué
  un círculo}: si el mes fuera el número 1 a 12, el modelo creería que diciembre (12) está lejos de enero (1); en el
  círculo quedan juntos;
\item si es uno de los cinco lugares (1 o 0).
\end{itemize}
Datos: 468 casos para el modelo final (234 tranquilos, 185 concurridos, 49 saturados).

\subsection{Los modelos y su fundamento}
\textbf{Línea base, persistencia}: $\hat y_{t+1}=y_t$ (``igual que este mes''). Un modelo que no le gane a esto no sirve.

\textbf{Regresión logística multiclase.} Para cada clase $k$ calcula un puntaje lineal y los convierte en
probabilidades con la función \emph{softmax}:
\[
z_k=b_k+\bm\beta_k^{\top}\bm x,\qquad P(y=k\mid \bm x)=\frac{e^{z_k}}{\sum_j e^{z_j}} .
\]
\emph{Fundamento}: el softmax garantiza probabilidades positivas que suman 1, y cada $\beta$ dice cuánto empuja un rasgo
hacia una clase. Los $\beta$ se estiman por \textbf{máxima verosimilitud}: se buscan los que hacen más probables los
estados que de verdad ocurrieron (equivale a minimizar la entropía cruzada). Detalles: los rasgos se \emph{estandarizan}
($(x-\bar x)/\sigma$) para que los $\beta$ sean comparables, y se usan \emph{pesos de clase balanceados}, porque
``saturado'' es raro (49 de 468) y sin pesos el modelo casi nunca lo predeciría.

\textbf{Random Forest}: 400 árboles de decisión; cada uno se entrena con una muestra al azar de los casos (con
reemplazo) y mira un subconjunto al azar de rasgos en cada corte. La predicción promedia las probabilidades de los 400.
\emph{Fundamento}: promediar muchos árboles distintos reduce la varianza de un árbol solo.

\textbf{Gradient Boosting}: suma árboles pequeños, $F_M(\bm x)=\sum_{m=1}^{M}\nu\,h_m(\bm x)$; cada árbol nuevo $h_m$ se
ajusta a los errores que dejaron los anteriores (al gradiente de la pérdida). 100 árboles con tasa $\nu=0.1$.

\begin{amano}[ · Mahahual, predicción de agosto de 2026 con la regresión logística]
\textbf{1. Estandarizar.} El índice de julio fue 0.731; la media de entrenamiento es 0.2762 y la desviación 0.2799:
\[x^{\text{std}}=\frac{0.731-0.2762}{0.2799}=1.625 .\]
Así con los siete rasgos: $(1.625,\ 1.362,\ 1.623,\ 0.871,\ -0.803,\ -1.158,\ -0.657)$.

\textbf{2. Puntaje de ``saturado''.} Coeficientes $\beta=(1.462,\ 1.291,\ 1.560,\ 0.567,\ -0.158,\ 1.195,\ -0.016)$ y
$b=-3.255$:
\begin{align*}
z_{\text{sat}}&=-3.255+1.462(1.625)+1.291(1.362)+1.560(1.623)+0.567(0.871)\\
&\quad+(-0.158)(-0.803)+1.195(-1.158)+(-0.016)(-0.657)\\
&=-3.255+2.376+1.758+2.532+0.494+0.127-1.384+0.011=2.659 .
\end{align*}
Igual se obtienen $z_{\text{conc}}=3.068$ y $z_{\text{tranq}}=-5.728$.

\textbf{3. Softmax.} Restando el mayor (3.068) para no desbordar:
\[
e^{0}=1,\quad e^{2.659-3.068}=e^{-0.409}=0.664,\quad e^{-5.728-3.068}=e^{-8.796}=0.00015 ;
\]
\[
P(\text{concurrido})=\frac{1}{1.664}=0.601,\qquad P(\text{saturado})=\frac{0.664}{1.664}=0.399,\qquad P(\text{tranquilo})\approx0 .
\]
El modelo dice: Mahahual sigue concurrido, con 40\,\% de probabilidad de saturarse. Es lo que muestra la página.
\end{amano}

\textbf{Qué rasgo pesa más} (promedio de $|\beta|$ en las tres clases): el índice del mes (1.07), el de dos meses antes
(1.04), el del mes anterior (0.95), la temporada como coseno (0.80), el cambio (0.40), ser uno de los cinco (0.35).

\subsection{Cómo se validó: origen móvil}
\textbf{Por qué no revolver al azar.} La forma escolar de validar (separar 80\,\% para entrenar y 20\,\% al azar para
probar) aquí sería trampa: el modelo podría entrenar con agosto y probarse con julio, es decir, \textbf{ver el futuro}. En
series de tiempo se valida como se usaría de verdad.

\textbf{Origen móvil}: para cada uno de los 12 meses de prueba (agosto de 2025 a julio de 2026), se reentrena cada modelo
con todo lo anterior y se predice solo ese mes. Son 12 reentrenamientos, 156 predicciones y 30 cambios reales de estado.

\subsection{Las métricas}
Para cada clase: precisión $P_k=\frac{VP_k}{VP_k+FP_k}$ (de lo que predijo $k$, cuánto era $k$), sensibilidad
$R_k=\frac{VP_k}{VP_k+FN_k}$ (de lo que era $k$, cuánto encontró) y $F1_k=\frac{2P_kR_k}{P_k+R_k}$ (su media armónica, que
castiga si una de las dos es baja). $F1_{\text{macro}}=\frac13\sum_k F1_k$ da el mismo peso a las tres clases, aunque
``saturado'' sea rara. Además: \textbf{cambios anticipados} (aciertos en los meses en que el estado cambió), que es lo que
le importa a la regla anti-colapso.

\begin{amano}[ · matriz de confusión y F1 de la logística]
\begin{center}
\begin{tabular}{lrrr}
\encabezado \h{Real $\backslash$ predicho} & \h{tranquilo} & \h{concurrido} & \h{saturado} \\
\textbf{tranquilo} & 75 & 12 & 0 \\ \textbf{concurrido} & 10 & 50 & 2 \\ \textbf{saturado} & 0 & 2 & 5 \\
\end{tabular}
\end{center}
\begin{itemize}
\item Tranquilo: $P=75/85=0.882$, $R=75/87=0.862$, $F1=\frac{2(0.882)(0.862)}{0.882+0.862}=0.872$.
\item Concurrido: $P=50/64=0.781$, $R=50/62=0.806$, $F1=0.794$. Saturado: $P=R=5/7=0.714$, $F1=0.714$.
\item $F1_{\text{macro}}=(0.872+0.794+0.714)/3=0.793$. Aciertos: $75+50+5=130$ de 156.
\end{itemize}
\end{amano}

\begin{center}
\begin{tabular}{lrrrr}
\encabezado \h{Modelo} & \h{Aciertos} & \h{Cambios (de 30)} & \h{Falsas alarmas} & \h{F1 macro} \\
Persistencia & 126 & 0 & 0 & 0.774 \\
\textbf{Regresión logística (elegida)} & \textbf{130} & \textbf{8} & 4 & 0.793 \\
Random Forest & 130 & 6 & 2 & 0.766 \\
Gradient Boosting & 128 & 7 & 5 & 0.731 \\
\end{tabular}
\end{center}

\begin{decision}
El modelo que más acierta y, a igualdad, el que más cambios anticipa (la regla anti-colapso necesita anticipar). Antes de
la auditoría ese criterio daba Random Forest; al agregar ``llegadas por cuarto'', la logística empató en aciertos y
anticipó más cambios. Se respeta el criterio, no el nombre del modelo: el código elige solo.
\end{decision}

\begin{supuestos}
\begin{itemize}
\item \textbf{Lectura honesta}: el estado se repite el 81\,\% de los meses, así que el modelo agrega poco (+4 aciertos);
  su valor está en los 8 cambios que anticipa.
\item \textbf{Sesgo}: en los cinco lugares hubo solo 2 cambios en 48 casos y el modelo no anticipó ninguno. La afirmación
  correcta es ``anticipa cambios en los lugares de referencia''; en el sur casi no hubo cambios con qué probarlo.
\end{itemize}
\end{supuestos}

\begin{codigoen}\codigo{radar/prediccion.py}: \codigo{tabla_de_aprendizaje}, \codigo{modelos}, \codigo{origen_movil},
\codigo{sesgo}, \codigo{predecir_mes_siguiente}.\end{codigoen}

\section{Pronóstico de visitantes: series de tiempo}
\label{sec:pronostico}
\subsection{De dónde salen las series}
\begin{center}
\begin{tabularx}{\linewidth}{lYrr}
\encabezado \h{Serie} & \h{Fuente y camino} & \h{Meses} & \h{Entrenan} \\
Visitantes a Oxtankah (Bahía) & DataTur \codigo{BdINAH.zip} → \codigo{silver/inah} → \codigo{pronostico/series.py} & 127 & 90 \\
Kohunlich + Dzibanché + Ichkabal (Ruta) & DataTur \codigo{BdINAH.zip} → \codigo{silver/inah} & 127 & 91 \\
Cruces desde Belice (Chetumal) & API de SITUR-Q → \codigo{silver/siturq} & 90 & 62 \\
Ocupación de Cancún (referencia) & DataTur semanal → \codigo{silver/datatur_ocupacion} & 55 & 55 \\
\bottomrule
\end{tabularx}
\end{center}
\textbf{Por qué estas}: el sur no tiene ocupación oficial en 2025–2026; se pronostica solo lo \textbf{medido} que llega a
2026. Maya Ka'an y la Laguna no tienen serie mensual propia: hueco declarado.

\subsection{Meses que no entrenan}
Un mes cerrado no es ``cero demanda''. Si el modelo aprendiera de esos ceros, creería que en abril puede haber 0
visitantes por temporada. Reglas (decisión ``Hueco + forma del año''):
\begin{itemize}
\item \textbf{Cierre}: la zona reportó 0 visitantes.
\item \textbf{Mes parcial}: el mes justo antes de cerrar o después de reabrir (abrió a medias).
\item \textbf{Pandemia}: INAH, de marzo de 2020 a diciembre de 2021; Belice, de marzo de 2020 a junio de 2022 (de marzo a
  junio de 2022 los cruces apenas iban en 65–82\,\% de 2019).
\item Una región toma el motivo más fuerte de sus zonas. Ejemplo, Ruta: diciembre de 2024 cierre; enero de 2025 cierre
  (Dzibanché seguía cerrada aunque Kohunlich abrió); febrero de 2025 parcial; marzo de 2025 ya entrena.
\end{itemize}
Esos meses no entrenan, pero su valor se conserva. Cada reapertura inicia un \textbf{tramo} con su propio nivel.

\subsection{Los cinco modelos}
Notación: $y_t$ = valor del mes $t$; $o$ = origen (último mes conocido); $h$ = meses hacia adelante; $S_m$ = forma del año
calculada \textbf{solo con años anteriores al origen}.

\paragraph{1. Línea base (ingenuo estacional).} $\hat y_{o+h}$ = el último valor del mismo mes del año. Es la vara.

\paragraph{2 y 3. Holt-Winters con la forma del año fija.}
\begin{intuicion}
Se le quita la temporada a la serie (se divide cada mes entre su índice), así queda el ``nivel'' de fondo. Ese nivel se
actualiza cada mes mezclando el dato nuevo con lo que ya se creía, y para pronosticar se le vuelve a poner la temporada.
\end{intuicion}
\[
z_t=\frac{y_t}{S_{m(t)}},\qquad \ell_t=\alpha z_t+(1-\alpha)\ell_{t-1},\qquad \hat y_{o+h}=\ell_o\,S_{m(o+h)} .
\]
$\alpha\in(0,1)$ dice cuánto se cree en el dato nuevo: cerca de 1, el nivel sigue cada dato; cerca de 0, cambia lento. Se
estima minimizando el error de pronóstico a un paso dentro del tramo (statsmodels). La variante \textbf{con tendencia
amortiguada} agrega una pendiente $b_t=\beta(\ell_t-\ell_{t-1})+(1-\beta)\phi b_{t-1}$ que se va apagando ($\phi<1$).
\emph{Por qué la forma fija}: el Holt-Winters clásico estima la temporada él mismo y necesita al menos 24 meses seguidos;
desde las reaperturas solo hay 17 a 20.

\begin{amano}[ · Ruta, origen julio de 2026, sin tendencia]
statsmodels estimó $\alpha=0.1183$. Julio tuvo 5,057 visitantes y su índice es 0.953: $z=5{,}057/0.953=5{,}306.6$. El
nivel de junio era 5,332.5:
\[
\ell_{\text{jul}}=0.1183\times5{,}306.6+0.8817\times5{,}332.5=627.8+4{,}701.7=5{,}329.4 .
\]
Enero de 2027: $\hat y=5{,}329.4\times S_1=5{,}329.4\times1.606=8{,}559$ visitantes. El código da 8,559.
\end{amano}

\paragraph{4. Regresión con clima} (mínimos cuadrados sobre el logaritmo):
\[
\ln y_t=\tau_{k(t)}+\mu_{m(t)}+\beta\,\frac{L_t-\bar L_{m(t)}}{100}+\gamma\,T_t+\varepsilon_t .
\]
$\tau_k$ = nivel del tramo $k$; $\mu_m$ = efecto del mes (enero es la referencia, $\mu_1=0$); $L_t-\bar L_m$ = lluvia del mes
menos la normal (en cientos de mm); $T_t=1$ si empezó una tormenta que afecta al sur.
\emph{Fundamento}: en logaritmos, los efectos se vuelven porcentajes: $e^{\mu_{12}}$ es ``cuántas veces enero'' es
diciembre. Los coeficientes se estiman por mínimos cuadrados ordinarios, $\hat{\bm\theta}=(X^\top X)^{-1}X^\top\bm y$
(numpy \codigo{lstsq}). Para el futuro el clima no se conoce, así que se usa el normal: lluvia anómala = 0 y tormenta =
su frecuencia histórica en ese mes:
\[
\hat y_{o+h}=\exp\big(\hat\tau_K+\hat\mu_m+\hat\gamma\,\hat p_m\big).
\]

\begin{amano}[ · Bahía, diciembre de 2026]
Con 90 meses que entrenan hay tres tramos: antes de 2020 ($\hat\tau=7.0975$, un enero típico de $e^{7.0975}=1{,}209$
visitantes), después de la pandemia ($6.8559$, 949) y desde la reapertura de 2024 ($7.0382$, 1,139). Diciembre:
$\hat\mu_{12}=0.1404$, o sea $e^{0.1404}=1.151$ veces un enero; nunca hubo tormenta en diciembre ($\hat p=0$):
\[
\hat y=e^{7.0382+0.1404}=e^{7.1786}=1{,}311\ \text{visitantes}.
\]
El código da 1,311.13. El coeficiente de la lluvia, $\hat\beta=-0.115$, dice que cada 100 mm sobre lo normal se asocian a
$e^{-0.115}-1=-10.9\,\%$ visitantes ($p=0.027$).
\end{amano}

\paragraph{5. Gradient Boosting con rezagos.} 150 árboles (tasa 0.05, hojas de al menos 20 casos) que aprenden de todos
los pares pasados (origen → destino), con: meses hacia adelante, mes del año, logaritmo del último valor, del promedio de
los últimos 3 meses y del mismo mes del año anterior.

\subsection{Cómo se midió el error: origen móvil}
Para cada mes del pasado (desde enero de 2019 en el INAH) se pronostican los 12 siguientes \textbf{usando solo lo que se
sabía entonces}, incluida la forma del año, y se compara con lo que pasó. Para que la comparación sea justa, solo cuentan
los pares que los cinco modelos pudieron pronosticar: 366 en la Bahía, 378 en la Ruta, 366 en Belice y 438 en Cancún.
\[
MAE=\frac1n\sum|y-\hat y|,\qquad MAPE=\frac{100}{n}\sum\left|\frac{y-\hat y}{y}\right|,\qquad
\text{error relativo}=\frac{MAE_{\text{modelo}}}{MAE_{\text{línea base}}} .
\]
El error relativo menor que 1 significa que el modelo le gana a ``el mismo mes del año pasado''.

\begin{amano}[ · el error de un pronóstico]
Bahía, origen enero de 2019, destino febrero: pronóstico 1,058.48, real 925.
$e=|\ln(1{,}058.48/925)|=|\ln 1.1443|=0.1348$: el pronóstico se pasó 14.4\,\%.
\end{amano}

\subsection{El rango del 90\,\%: intervalo conformal}
\begin{intuicion}
Se revisan los errores que el modelo cometió en el pasado en pronósticos parecidos. Si 9 de cada 10 fueron menores que
cierto tamaño, lo razonable es decir que el próximo pronóstico tampoco se equivocará más que eso, con 90\,\% de confianza.
No se supone que los errores sigan una curva normal: se usan los errores reales.
\end{intuicion}
\[
\hat q=e_{(k)},\quad k=\left\lceil (n+1)\cdot0.9\right\rceil,\qquad \text{rango}=\left[\hat y\,e^{-\hat q},\ \hat y\,e^{\hat q}\right] .
\]
$e_{(k)}$ es el $k$-ésimo error más chico de los $n$ errores de calibración (del mismo modelo y del mismo tramo de
horizonte: 1–3, 4–6 o 7–12 meses) cuyo mes ya había pasado en el origen.

\textbf{Fundamento (la garantía)}: si los errores futuros son intercambiables con los pasados (vienen ``del mismo tipo de
mes''), el nuevo error cae por debajo del $k$-ésimo con probabilidad al menos $k/(n+1)$. El $+1$ y el techo son la
corrección de muestra finita que vuelve la garantía exacta. Se trabaja con errores en logaritmo porque los errores son
proporcionales al nivel.

\begin{amano}[ · Bahía, diciembre de 2026 (5 meses adelante)]
Hay $n=105$ errores de calibración: $k=\lceil 106\times0.9\rceil=\lceil 95.4\rceil=96$, y el 96.º más chico es
$\hat q=0.3132$. Garantía: $96/106=90.6\,\%$. Con el pronóstico de 1,311.13:
\begin{align*}
\text{mínimo}&=1{,}311.13\times e^{-0.3132}=1{,}311.13\times0.7311=959,\\
\text{máximo}&=1{,}311.13\times e^{0.3132}=1{,}311.13\times1.3678=1{,}793 .
\end{align*}
Diciembre de 2025 tuvo 1,224: dentro del rango.
\end{amano}

\textbf{Cobertura real}: se midió cuántas veces el dato cayó dentro del rango, y se reporta aunque salga peor.

\subsection{La elección y los resultados}
\begin{decision}
``Menor error con rango ≥ 80\,\%'': por serie, el modelo de menor error entre los que tienen un rango que se cumple al
menos 80 de cada 100 veces. Descartados: el menor error sin importar el rango (en la Ruta elegiría un modelo cuyo rango
acierta 72 de cada 100) y un solo modelo para todo (en Cancún la regresión comete 58\,\% más error que la línea base).
\end{decision}

\begin{center}
\small
\begin{tabular}{llrrrr}
\encabezado \h{Serie} & \h{Modelo} & \h{MAPE} & \h{vs base} & \h{Cobertura} & \h{Rango típico} \\
Bahía & \textbf{Regresión con clima} & 15.7\,\% & \textbf{0.898} & 91.3\,\% & $\pm$43\,\% \\
 & Holt-Winters sin tendencia & 19.4\,\% & 1.078 & 96.5\,\% & $\pm$63\,\% \\
 & Holt-Winters con tendencia & 27.6\,\% & 1.524 & 87.8\,\% & $\pm$108\,\% \\
 & Gradient Boosting & 28.7\,\% & 1.405 & 85.8\,\% & $\pm$83\,\% \\
 & Línea base & 19.8\,\% & 1.000 & 85.4\,\% & $\pm$77\,\% \\
Ruta & \textbf{Regresión con clima} & 21.6\,\% & \textbf{0.737} & 80.3\,\% & $\pm$51\,\% \\
 & Holt-Winters sin tendencia & 20.1\,\% & 0.639 & 72.3\,\% & $\pm$45\,\% \\
 & Línea base & 27.9\,\% & 1.000 & 74.3\,\% & $\pm$63\,\% \\
Belice & \textbf{Regresión con clima} & 11.0\,\% & \textbf{0.820} & 94.1\,\% & $\pm$26\,\% \\
Cancún & \textbf{Línea base} & 3.8\,\% & \textbf{1.000} & 89.3\,\% & $\pm$10\,\% \\
\end{tabular}
\normalsize
\end{center}

\begin{supuestos}
\begin{itemize}
\item \textbf{La Ruta no llega a 90\,\% con ningún modelo}: en 2023 las visitas cayeron 10.7\,\% (de 46,295 a 41,322) sin
  aviso en la historia, y ese año la cobertura fue de 67\,\%. La garantía conformal supone que el futuro se parece al
  pasado; 2023 rompió ese supuesto.
\item \textbf{El clima casi no mejora el pronóstico}: la misma regresión sin lluvia ni tormentas da errores casi iguales
  (0.906, 0.842 y 0.721). Lo que gana es la estructura: nivel por tramo y efecto del mes. Al pronosticar se usa el clima
  normal, porque el futuro no se conoce.
\item \textbf{La tendencia falla en tramos cortos}: Holt-Winters con tendencia llegó a pronosticar $-240$ visitantes en la
  Bahía porque aprendió la tendencia en plena reapertura.
\item \textbf{En Cancún, repetir el año anterior es lo mejor} (3.8\,\% de error): su ocupación cambia poco de un año a otro.
\end{itemize}
\end{supuestos}

\begin{figure}[h]\centering
\includegraphics[width=0.85\linewidth]{f11_pronostico_12_meses.png}
\caption{Pronóstico de los próximos 12 meses con su rango del 90\,\% (notebook 03).}
\end{figure}

\begin{codigoen}\codigo{pronostico/series.py}, \codigo{modelos.py} (cinco modelos y \codigo{origen_movil}),
\codigo{intervalos.py} (\codigo{cuantil_conformal}, \codigo{cobertura}), \codigo{seleccion.py}.\end{codigoen}

\begin{teprofe}
\emph{``¿Por qué no usaste ARIMA o Prophet?''} Porque las series tienen huecos largos (pandemia y cierres por obras) y
tramos cortos desde las reaperturas, y esos métodos necesitan series continuas largas. Se eligieron modelos que funcionan
por tramos y con la temporada estimada aparte, y se compararon contra una línea base con validación honesta.
\end{teprofe}
